\documentclass[a4paper]{article}

\usepackage[spanish]{babel}
\usepackage{graphicx}
\usepackage{amsmath, amssymb}
\usepackage[margin=2cm]{geometry}
\usepackage{fancyhdr}
\usepackage{enumerate}
\usepackage[shortlabels]{enumitem}
\usepackage{parskip}
\usepackage[most]{tcolorbox}
\usepackage[hidelinks]{hyperref}
\usepackage{float}
\usepackage{booktabs}



% cabecera
\pagestyle{fancy}
\fancyhead[l]{Astroestadística}
\fancyhead[c]{Problemset 1}
\fancyhead[r]{\today}
\fancyfoot[c]{\thepage}
\renewcommand{\headrulewidth}{0.2pt} % linea horizontal

\begin{document}

\begin{titlepage}
  \centering
  {\Large Universidad de Antioquia\par}
  \vspace{2cm}
  {\huge\bfseries Astroestadística\par}
  \vspace{0.4cm}
  {\Large Solución al problemset 1 \par}
  \vspace{2.5cm}
  {\large Autor:\par}
  \vspace{0.2cm}
  {\large Daniel Soto Villada\par}
  {\large Estudiante de Astronomía\par}
  \vfill
  {\large \today\par}
\end{titlepage}

\newpage
\tableofcontents
\newpage


\section{Problema 1}
\textbf{Enunciado:}
Estrellas. Suponga que las estrellas en un cúmulo estelar se pueden clasificar según 2 características: color (\textbf{rojas} y \textbf{azules}) y tamaño (\textbf{gigantes} y \textbf{enanas}). Un experimento consiste en observar una estrella de dicho cúmulo de forma aleatoria, considerando que el 80\% de las estrellas son rojas, el 70\% de las estrellas rojas son enanas y el 90\% de las estrellas azules son gigantes. Sea $A$ el suceso o evento en el cual se observa una estrella roja y $B$ el suceso en el cual se observa una estrella gigante. Describa los siguientes sucesos y calcule su probabilidad basándose en los axiomas y propiedades de la probabilidad:
\begin{enumerate}[a.]
    \item $A \cup B$
    \item $A \cap B$
    \item $A \cup B'$
    \item $A' \cup B'$
    \item $A - B$
    \item $A' - B'$
    \item $(A \cap B) \cup (A \cap B')$
\end{enumerate}

\subsection*{Solución}

\textbf{Interpretación en notación matemática:}

Para resolver este problema, primero debemos definir los eventos fundamentales y traducir los porcentajes proporcionados a valores de probabilidad según la teoría de Devore. Definimos el espacio muestral $S$ como el conjunto de todas las estrellas del cúmulo.

Eventos principales:
\begin{itemize}
    \item $A$: La estrella es roja.
    \item $A'$: La estrella es azul (complemento de $A$, dado que solo hay dos colores).
    \item $B$: La estrella es gigante.
    \item $B'$: La estrella es enana (complemento de $B$).
\end{itemize}

\textbf{Información proporcionada:}
\begin{itemize}
    \item $P(A) = 0.80$ (Proporción de estrellas rojas).
    \item $P(B' | A) = 0.70$ (Probabilidad condicional de ser enana dado que es roja).
    \item $P(B | A') = 0.90$ (Probabilidad condicional de ser gigante dado que es azul).
\end{itemize}

\textbf{Información a hallar:}
Se requiere describir conceptualmente y calcular numéricamente la probabilidad de diversas operaciones entre estos conjuntos astronómicos.

\textbf{Cálculos preliminares:}
Antes de proceder con los incisos, calculamos las probabilidades de las intersecciones y las marginales necesarias usando la regla de la multiplicación y la ley de probabilidad total:
\begin{enumerate}
    \item $P(A') = 1 - P(A) = 1 - 0.80 = 0.20$.
    \item $P(B | A) = 1 - P(B' | A) = 1 - 0.70 = 0.30$.
    \item $P(A \cap B) = P(B | A) \cdot P(A) = (0.30)(0.80) = 0.24$.
    \item $P(A \cap B') = P(B' | A) \cdot P(A) = (0.70)(0.80) = 0.56$.
    \item $P(A' \cap B) = P(B | A') \cdot P(A') = (0.90)(0.20) = 0.18$.
    \item $P(A' \cap B') = P(B' | A') \cdot P(A') = (1 - 0.90)(0.20) = (0.10)(0.20) = 0.02$.
    \item $P(B) = P(A \cap B) + P(A' \cap B) = 0.24 + 0.18 = 0.42$.
    \item $P(B') = 1 - P(B) = 0.58$.
\end{enumerate}

\textbf{Solución de los incisos:}

\begin{enumerate}[a.]
    \item \textbf{$A \cup B$ (Unión):} Representa el evento en que la estrella es \textit{roja o gigante} (incluyendo las que son ambas). Según la regla de la adición:
    $$P(A \cup B) = P(A) + P(B) - P(A \cap B) = 0.80 + 0.42 - 0.24 = 0.98$$

    \item \textbf{$A \cap B$ (Intersección):} Es el evento en que la estrella es \textit{roja y gigante} simultáneamente. Según el cálculo previo:
    $$P(A \cap B) = 0.24$$

    \item \textbf{$A \cup B'$ (Unión con el complemento):} Representa el evento en que la estrella es \textit{roja o enana}. Usando la regla de adición:
    $$P(A \cup B') = P(A) + P(B') - P(A \cap B') = 0.80 + 0.58 - 0.56 = 0.82$$

    \item \textbf{$A' \cup B'$ (Unión de complementos):} Representa que la estrella \textit{no es roja o no es gigante}. Por las leyes de De Morgan, esto equivale al complemento de la intersección $(A \cap B)'$:
    $$P(A' \cup B') = 1 - P(A \cap B) = 1 - 0.24 = 0.76$$

    \item \textbf{$A - B$ (Diferencia):} Representa el evento de observar una estrella que es \textit{roja pero no gigante} (es decir, una estrella roja enana). Es equivalente a $A \cap B'$:
    $$P(A - B) = P(A) - P(A \cap B) = 0.80 - 0.24 = 0.56$$

    \item \textbf{$A' - B'$ (Diferencia de complementos):} Describe una estrella que \textit{no es roja y no es enana}. Esto equivale a observar una estrella azul gigante ($A' \cap B$):
    $$P(A' - B') = P(A') - P(A' \cap B') = 0.20 - 0.02 = 0.18$$

    \item \textbf{$(A \cap B) \cup (A \cap B')$ (Propiedad distributiva):} Este evento describe una estrella que es \textit{(roja y gigante) o (roja y enana)}. Dado que ser gigante y enana son eventos mutuamente excluyentes, esta unión recompone exactamente el evento $A$ (la estrella es roja). Estadísticamente, estamos sumando las particiones del evento $A$ basadas en $B$:
    $$P((A \cap B) \cup (A \cap B')) = P(A) = 0.80$$
\end{enumerate}


\section{Problema 2}
\textbf{Enunciado:}
2. Clasificación de galaxias. Basándose en el diagrama de la secuencia de Hubble (tuning fork) proporcionado en la imagen del conjunto de problemas:
\begin{enumerate}[a.]
    \item Identifique cuántas subcategorías existen para las galaxias elípticas (E) según su grado de elipticidad.
    \item Determine cuántas clasificaciones posibles existen para las galaxias espirales, considerando tanto las familias no barradas (S) como las barradas (SB) y sus respectivas subcategorías morfológicas (a, b, c).
    \item Calcule el número total de clasificaciones morfológicas distintas presentes en el diagrama completo.
    \item Si se selecciona una galaxia al azar de este esquema y se asume que cada clasificación es igualmente probable, ¿cuál es la probabilidad de que la galaxia sea de tipo espiral barrada?
\end{enumerate}

\subsection*{Solución}

\textbf{Interpretación en notación matemática:}
Este problema se resuelve mediante técnicas de conteo, específicamente la **regla del producto** para experimentos en etapas y la enumeración de elementos en un espacio muestral discreto $S$.

\begin{itemize}
    \item Sea $n_E$ el número de categorías elípticas.
    \item Sea $n_{S0}$ el número de categorías lenticulares.
    \item Sea $n_S$ el número de categorías espirales no barradas.
    \item Sea $n_{SB}$ el número de categorías espirales barradas.
\end{itemize}

\textbf{Información proporcionada (según la tabla/diagrama visual):}
\begin{itemize}
    \item Elípticas: El índice numérico va de 0 a 7.
    \item Lenticulares: Representadas por S0.
    \item Espirales: Dos ramas con tres subcategorías cada una (a, b, c).
\end{itemize}

\textbf{Información a hallar:}
El conteo total y las probabilidades asociadas bajo el supuesto de equiprobabilidad.

\textbf{Desarrollo extenso del problema:}

\begin{enumerate}[a.]
    \item \textbf{Conteo de galaxias elípticas:} 
    Observando la información numérica del eje horizontal del diagrama, las galaxias elípticas se clasifican desde E0 hasta E7. Contando los elementos:
    $$ n_E = \{0, 1, 2, 3, 4, 5, 6, 7\} \implies n_E = 8 $$
    Existen **8** formas de clasificar una galaxia elíptica según este esquema.

    \item \textbf{Conteo de galaxias espirales:}
    Para las galaxias espirales, el experimento tiene dos etapas: 
    1. Selección de la familia (Barrada $SB$ o No barrada $S$): $n_{fam} = 2$.
    2. Selección de la apertura de brazos ($a, b, c$): $n_{sub} = 3$.
    Aplicando la **regla del producto** (Devore, Sección 2.3):
    $$ N_{espirales} = n_{fam} \cdot n_{sub} = 2 \cdot 3 = 6 $$
    Las clasificaciones son: Sa, Sb, Sc, SBa, SBb, SBc.

    \item \textbf{Número total de clasificaciones ($N$):}
    El total es la suma de todas las categorías mutuamente excluyentes presentes en el diagrama (Elípticas + Lenticulares + Espirales):
    $$ N = n_E + n_{S0} + n_S + n_{SB} $$
    $$ N = 8 + 1 + 3 + 3 = 15 $$
    El diagrama muestra un total de **15** clasificaciones morfológicas distintas.

    \item \textbf{Cálculo de probabilidad:}
    Si los resultados son igualmente probables, la probabilidad de un evento $A$ es $P(A) = \frac{N(A)}{N}$. Sea $A$ el evento de observar una espiral barrada ($SB$):
    - Resultados en $A$: {SBa, SBb, SBc} $\implies N(A) = 3$.
    - Resultados totales: $N = 15$.
    $$ P(SB) = \frac{3}{15} = \frac{1}{5} = 0.20 $$
    La probabilidad de observar una galaxia espiral barrada es del $20\%$.
\end{enumerate}

\section{Problema 3}
\textbf{Enunciado:}
3. Sistemas binarios. La mayoŕıa de las estrellas que vemos en el cielo son sistemas binarios, con dos (o a veces más) estrellas orbitando un centro de masa común. Tal vez el 85\% de las estrellas son sistemas binarios. Además, se estima que el 50–60\% de las estrellas binarias son capaces de soportar planetas tipo terrestres habitables dentro de los rangos orbitales estables.
\begin{enumerate}[a.]
    \item ¿Cuál es la probabilidad de que, al mirar 10 estrellas en una noche de forma aleatoria, todas sean sistemas binarios? Suponga que el universo es homogéneo e isotrópico.
    \item ¿Cuál es la probabilidad de que, al mirar una estrella de forma aleatoria, se esté mirando un sistema binario capaz de soportar planetas tipo terrestres habitables?
\end{enumerate}

\subsection*{Solución}

\textbf{Interpretación en notación matemática:}

Este problema se aborda utilizando la teoría de \textbf{probabilidad elemental} y la \textbf{distribución binomial}, fundamentadas en el texto de Devore. Definimos los sucesos o eventos relevantes en nuestro espacio muestral estelar:

\begin{itemize}
    \item \textbf{Evento $B$:} La estrella observada pertenece a un sistema binario.
    \item \textbf{Evento $H$:} El sistema estelar es capaz de soportar planetas tipo terrestres habitables.
\end{itemize}

\textbf{Información proporcionada:}
\begin{itemize}
    \item $P(B) = 0.85$: Probabilidad de que una estrella sea un sistema binario (85\% de la población estelar).
    \item $P(H|B) \in [0.50, 0.60]$: Probabilidad condicional de que un sistema soporte planetas habitables, dado que es binario (entre el 50\% y el 60\%).
    \item $n = 10$: Tamaño de la muestra (estrellas observadas) para el inciso (a).
\end{itemize}

\textbf{Información a hallar:}
\begin{itemize}
    \item En el inciso \textbf{a}, se requiere la probabilidad de obtener 10 éxitos en 10 ensayos independientes bajo un modelo binomial.
    \item En el inciso \textbf{b}, se requiere la probabilidad de la intersección de dos eventos, $P(B \cap H)$.
\end{itemize}

\textbf{Desarrollo de la solución:}

\begin{enumerate}[a.]
    \item \textbf{Análisis de probabilidad conjunta para eventos independientes:}
    Para calcular la probabilidad de que 10 estrellas observadas sean todas binarias, debemos considerar cada observación como un \textit{ensayo de Bernoulli} independiente. Según el enunciado, la suposición de un universo \textit{homogéneo e isotrópico} garantiza que la probabilidad de éxito $p$ sea constante en cada observación y que los resultados sean independientes entre sí.
    
    Sea $X$ la variable aleatoria que cuenta el número de sistemas binarios en la muestra. Bajo estas condiciones, $X$ sigue una distribución binomial con parámetros $n=10$ y $p=0.85$. La función masa de probabilidad (fmp) es:
    $$ b(x; n, p) = \binom{n}{x} p^x (1-p)^{n-x} $$
    
    Para hallar la probabilidad de que \textit{todas} sean sistemas binarios ($x=10$):
    $$ P(X = 10) = \binom{10}{10} (0.85)^{10} (0.15)^0 $$
    $$ P(X = 10) = 1 \cdot (0.85)^{10} \cdot 1 $$
    $$ P(X = 10) \approx 0.19687 $$
    
    Por lo tanto, la probabilidad de que las 10 estrellas sean sistemas binarios es aproximadamente del \textbf{19.69\%}.

    \item \textbf{Cálculo de la probabilidad de la intersección:}
    Se nos pide la probabilidad de observar una estrella que sea \textit{simultáneamente} un sistema binario \textbf{y} capaz de soportar planetas habitables. Según la regla de multiplicación de probabilidades (Devore, Sección 2.4):
    $$ P(B \cap H) = P(H|B) \cdot P(B) $$
    
    Dado que la probabilidad condicional $P(H|B)$ se proporciona como un rango ($50-60\%$), el resultado se puede expresar mediante el valor esperado promedio del rango (55\%) o como un intervalo de probabilidad:
    \begin{itemize}
        \item \textbf{Límite inferior:} $0.50 \cdot 0.85 = 0.425$
        \item \textbf{Límite superior:} $0.60 \cdot 0.85 = 0.510$
    \end{itemize}
    
    Si tomamos el valor medio del estimado ($55\%$):
    $$ P(B \cap H) = 0.55 \cdot 0.85 = 0.4675 $$
    
    Estadísticamente, esto se interpreta como que existe aproximadamente un \textbf{46.75\%} de probabilidad de que una estrella seleccionada al azar cumpla con ambas condiciones astronómicas.
\end{enumerate}

\section{Problema 4}
\textbf{Enunciado:}
4. La bolita roja. La caja I contiene 3 bolas rojas y 7 azules, la caja II contiene 8 bolas rojas y 2 azules, y la caja III contiene 4 bolas rojas y 6 azules. Se requiere sacar la bola roja para ganarse un regalito.
\begin{enumerate}[a.]
    \item ¿Con cuál caja se tiene mayor probabilidad de ganar?
    \item Usted no conoce el contenido de cada caja, así que debe elegir una al azar. ¿Qué probabilidad tiene de ganar?
    \item Suponga que usted es un feliz ganador, ¿cuál es la probabilidad de haber sacado la bolita de la caja III? ¿Y de la caja II?
\end{enumerate}

\subsection*{Solución}

\textbf{Interpretación en notación matemática:}

Para resolver este problema, aplicamos los conceptos de \textbf{probabilidad condicional}, la \textbf{Ley de Probabilidad Total} y el \textbf{Teorema de Bayes}, siguiendo la teoría expuesta en el texto de Devore. Primero, definimos los eventos relevantes:
\begin{itemize}
    \item $C_I, C_{II}, C_{III}$: Eventos en los que se selecciona la Caja I, II o III, respectivamente.
    \item $R$: Evento en el que se extrae una bola roja (ganar).
    \item $A$: Evento en el que se extrae una bola azul (perder).
\end{itemize}

\textbf{Información proporcionada:}
Basándonos en la composición de cada caja (donde cada una tiene un total de 10 bolas), las probabilidades condicionales de éxito dado que se elige una caja específica son:
\begin{itemize}
    \item $P(R|C_I) = \frac{3}{10} = 0.30$ (o $30\%$)
    \item $P(R|C_{II}) = \frac{8}{10} = 0.80$ (o $80\%$)
    \item $P(R|C_{III}) = \frac{4}{10} = 0.40$ (o $40\%$)
\end{itemize}
Para el inciso b, al elegir una caja al azar sin conocer su contenido, asignamos probabilidades iguales a la selección de cada caja (\textit{priors} equiprobables):
\begin{itemize}
    \item $P(C_I) = P(C_{II}) = P(C_{III}) = \frac{1}{3} \approx 0.3333$
\end{itemize}

\textbf{Información a hallar:}
\begin{itemize}
    \item Identificar el máximo entre las probabilidades condicionales.
    \item La probabilidad marginal de ganar, $P(R)$.
    \item Las probabilidades posteriores $P(C_{III}|R)$ y $P(C_{II}|R)$.
\end{itemize}

\textbf{Desarrollo extenso del problema:}

\begin{enumerate}[a.]
    \item \textbf{Identificación de la mejor opción:}
    Para determinar con cuál caja se tiene mayor probabilidad de ganar, simplemente comparamos las probabilidades condicionales calculadas a partir de la composición de las cajas.
    $$ P(R|C_{II}) = 0.80 > P(R|C_{III}) = 0.40 > P(R|C_I) = 0.30 $$
    Se tiene la mayor probabilidad de ganar con la \textbf{Caja II}, ya que posee la mayor proporción de bolas rojas en su interior.

    \item \textbf{Cálculo de la probabilidad total:}
    Cuando el experimento se realiza en dos etapas (primero elegir una caja al azar y luego extraer la bola), el evento de interés es la probabilidad marginal $P(R)$. Según la \textbf{Ley de Probabilidad Total} (Devore, Sección 2.4):
    $$ P(R) = P(R|C_I)P(C_I) + P(R|C_{II})P(C_{II}) + P(R|C_{III})P(C_{III}) $$
    Sustituyendo los valores:
    $$ P(R) = (0.30)(\frac{1}{3}) + (0.80)(\frac{1}{3}) + (0.40)(\frac{1}{3}) $$
    $$ P(R) = \frac{0.30 + 0.80 + 0.40}{3} = \frac{1.50}{3} = 0.50 $$
    La probabilidad de ganar bajo estas condiciones es del \textbf{50\%}.

    \item \textbf{Aplicación del Teorema de Bayes:}
    Si sabemos que el evento $R$ ya ocurrió (el jugador ganó), queremos actualizar nuestra creencia sobre de qué caja provino la bola. Utilizamos el \textbf{Teorema de Bayes} (Devore, Sección 2.4), que relaciona las probabilidades posteriores con la verosimilitud y los priors:
    $$ P(C_j|R) = \frac{P(R|C_j)P(C_j)}{P(R)} $$
    
    \textbf{Probabilidad para la Caja III:}
    $$ P(C_{III}|R) = \frac{(0.40)(\frac{1}{3})}{0.50} = \frac{0.1333}{0.50} = 0.2666... \approx 26.67\% $$
    (Equivalentemente, $\frac{4/30}{1/2} = \frac{8}{30} = \frac{4}{15}$).

    \textbf{Probabilidad para la Caja II:}
    $$ P(C_{II}|R) = \frac{(0.80)(\frac{1}{3})}{0.50} = \frac{0.2666}{0.50} = 0.5333... \approx 53.33\% $$
    (Equivalentemente, $\frac{8/30}{1/2} = \frac{16}{30} = \frac{8}{15}$).
\end{enumerate}



\section{Problema 5}
\textbf{Enunciado:}
5. Combinaciones con repetición. Una combinación con repetición de orden $r$ de los $n$ elementos de $A$ es una selección no ordenada de $r$ elementos de $A$ que pueden repetirse. El número de tales combinaciones se denota $CR_{n,r}$ y está dado por:
$$ CR_{n,r} = \binom{n + r - 1}{r} $$
\begin{enumerate}[a.]
    \item Demuestre la expresión anterior (Ayuda: Consulte sobre el método de barras y estrellas).
    \item De cuántas formas se pueden elegir 15 estrellas entre 7 tipos espectrales si:
    \begin{itemize}
        \item Se pueden elegir de cualquier forma.
        \item Se debe elegir una de cada tipo.
    \end{itemize}
\end{enumerate}

\subsection*{Solución}

\textbf{Interpretación en notación matemática:}

Este problema se fundamenta en las \textbf{técnicas de conteo} dentro de la teoría de la probabilidad. Según Devore, cuando el orden en que se seleccionan los elementos no es importante, pero se permite que un mismo elemento sea seleccionado más de una vez, estamos ante un escenario de \textbf{combinaciones con repetición}. 

Matemáticamente, si tenemos un conjunto con $n$ categorías distintas (en este caso, tipos espectrales) y deseamos realizar $r$ extracciones (selección de 15 estrellas), el problema consiste en hallar el número de formas de asignar $r$ objetos idénticos en $n$ contenedores distinguibles.

\textbf{Información proporcionada:}
\begin{itemize}
    \item Definición de combinación con repetición.
    \item Fórmula general: $CR_{n,r} = \binom{n+r-1}{r}$.
    \item Para el inciso (b): $n = 7$ tipos espectrales y $r = 15$ estrellas a elegir.
\end{itemize}

\textbf{Información a hallar:}
\begin{itemize}
    \item Demostración lógica de la fórmula mediante el método de barras y estrellas.
    \item Cálculo numérico del total de formas de elección bajo dos restricciones distintas.
\end{itemize}

\textbf{Desarrollo extenso del problema:}

\begin{enumerate}[a.]
    \item \textbf{Demostración mediante el método de barras y estrellas:}
    
    Para visualizar este problema, imaginemos que representamos cada uno de los $r$ elementos seleccionados como \textbf{estrellas} ($\star$). Para distinguir entre las $n$ categorías o tipos, utilizamos \textbf{barras} ($|$) como separadores.
    
    Si tenemos $n$ tipos de elementos, necesitaremos exactamente $n - 1$ barras para crear los compartimentos necesarios. Por ejemplo, si $n=3$, las barras $|$ $|$ dividen el espacio en tres secciones (Tipo 1, Tipo 2 y Tipo 3). Una selección de 5 elementos ($r=5$) podría verse así: $\star \star | \star | \star \star$. Esto indicaría 2 elementos del tipo 1, 1 del tipo 2 y 2 del tipo 3.
    
    El problema de conteo se reduce entonces a determinar cuántas formas hay de disponer un total de $r + (n - 1)$ símbolos en una línea, donde $r$ son estrellas y $n - 1$ son barras. Según la teoría de combinaciones de Devore, el número de formas de elegir $k$ posiciones para un símbolo (estrellas) de un total de $m$ posiciones disponibles es $\binom{m}{k}$.
    
    Sustituyendo nuestros valores:
    \begin{itemize}
        \item Total de posiciones ($m$): $r + (n - 1)$.
        \item Elementos a posicionar ($k$): $r$ (estrellas).
    \end{itemize}
    
    Por lo tanto, el número de combinaciones únicas es:
    $$ CR_{n,r} = \binom{n + r - 1}{r} $$
    Lo cual demuestra la expresión dada en el enunciado.

    \item \textbf{Cálculo para la elección de estrellas:}
    
    En este escenario astronómico, tenemos $n = 7$ (tipos espectrales) y $r = 15$ (estrellas totales a elegir).
    
    \begin{itemize}
        \item \textbf{Caso 1: Elección de cualquier forma.}
        Aquí aplicamos la fórmula directamente, ya que no hay restricciones sobre cuántas estrellas de cada tipo deben elegirse (puede haber tipos con cero estrellas).
        $$ n = 7, r = 15 $$
        $$ CR_{7,15} = \binom{7 + 15 - 1}{15} = \binom{21}{15} $$
        Usando la propiedad de simetría de las combinaciones $\binom{n}{k} = \binom{n}{n-k}$ mencionada en las fuentes:
        $$ \binom{21}{15} = \binom{21}{6} = \frac{21 \cdot 20 \cdot 19 \cdot 18 \cdot 17 \cdot 16}{6 \cdot 5 \cdot 4 \cdot 3 \cdot 2 \cdot 1} = 54,264 $$
        Existen \textbf{54,264} formas de elegir las estrellas.

        \item \textbf{Caso 2: Se debe elegir al menos una de cada tipo.}
        Si es obligatorio tener una estrella de cada uno de los 7 tipos, primero "apartamos" o asignamos 7 estrellas (una para cada tipo). Esto nos deja con $15 - 7 = 8$ estrellas restantes por elegir libremente entre los 7 tipos.
        
        Ahora el problema se transforma en elegir $r' = 8$ estrellas de $n = 7$ tipos con repetición:
        $$ CR_{7,8} = \binom{7 + 8 - 1}{8} = \binom{14}{8} $$
        Calculando el valor:
        $$ \binom{14}{8} = \binom{14}{6} = \frac{14 \cdot 13 \cdot 12 \cdot 11 \cdot 10 \cdot 9}{6 \cdot 5 \cdot 4 \cdot 3 \cdot 2 \cdot 1} = 3,003 $$
        Existen \textbf{3,003} formas de realizar la elección asegurando la presencia de todos los tipos espectrales.
    \end{itemize}
\end{enumerate}






\bibliographystyle{plainurl}
\bibliography{bibliografia} 
\end{document}
